\section{Compute-Optimal Test-Time Scaling}

\subsection{并行采样与串行修订}

增加 test-time compute 有两种基本方式：

\begin{itemize}
    \item \textbf{Parallel sampling}：从同一 prompt 独立生成多个候选，强调多样性；
    \item \textbf{Sequential revision}：让模型读取中间结果并修订，强调单轨迹深度。
\end{itemize}

\videofigure{figures/parallel-sampling.jpg}{并行采样用多个独立候选扩大搜索宽度。}{00:27:00--00:29:06}

并行方法易于扩展，且错误不会跨候选传播；串行方法能利用反馈逐步改进，但若早期方向错误，后续步骤可能围绕错误前提继续优化。

\subsection{Outcome Reward Model 与 Process Reward Model}

\videofigure{figures/reward-model-types.jpg}{Outcome reward 与 process reward 提供不同粒度的搜索信号。}{00:29:06--00:30:08}

\begin{itemize}
    \item \textbf{ORM}（Outcome Reward Model）：只评价最终答案；适合 best-of-$N$，信用分配粗；
    \item \textbf{PRM}（Process Reward Model）：逐步评价推理过程；可用于 beam search 和提前剪枝。
\end{itemize}

若一条轨迹包含步骤 $z_{1:T}$，PRM 可提供局部得分 $r_t=r(z_t\mid z_{<t},x)$，搜索器据此保留高潜力分支，而无需等到完整答案生成后才发现错误。

\subsection{混合宽度与深度}

Compute-optimal 策略不是固定选并行或串行，而是在总预算下组合两者：先探索多个初始候选，再对高分候选进行多轮修订。

\videofigure{figures/hybrid-search.jpg}{并行采样扩大宽度，revision model 增加搜索深度。}{00:31:10--00:35:28}

可写成预算约束：

$$
B=n_{\mathrm{parallel}}\cdot c_{\mathrm{sample}}
+n_{\mathrm{revise}}\cdot c_{\mathrm{revise}}
+c_{\mathrm{verify}}.
$$

\begin{itemize}
    \item $B$：单题总推理预算；
    \item $n_{\mathrm{parallel}}$：并行候选数；
    \item $n_{\mathrm{revise}}$：串行修订轮数；
    \item $c_{\mathrm{sample}},c_{\mathrm{revise}},c_{\mathrm{verify}}$：各操作成本。
\end{itemize}

\subsection{实验结论：最佳策略随难度变化}

\videofigure{figures/compute-optimal-results.jpg}{在 PRM 指导下，compute-optimal search 早期收益高于纯 best-of-$N$。}{00:35:28--00:37:28}

\videofigure{figures/parallel-sequential-ratio.jpg}{简单题偏向串行修订，困难题需要一定并行探索。}{00:37:28--00:40:12}

直觉是：

\begin{itemize}
    \item 简单题通常已有接近正确的初始解，revision 能快速修正；
    \item 困难题的初始解可能落入完全错误的思路，需要并行探索不同方向；
    \item 过度串行会把预算浪费在修补不可救的轨迹；
    \item 过度并行则产生大量浅层候选，没有充分利用反馈。
\end{itemize}

\subsection{测试时 FLOPs 何时胜过预训练 FLOPs}

\videofigure{figures/testtime-vs-pretraining.jpg}{中等难度问题常从 test-time compute 获益，而最难问题仍受基础能力限制。}{00:40:12--00:42:32}

\begin{importantbox}{Test-time scaling 不是 pre-training scaling 的完全替代}
当基础模型对某类问题几乎没有正确轨迹时，搜索无法放大不存在的信号。实验显示，对最难问题，增加推理计算的收益很小，扩大预训练或改进训练数据仍更有效。
\end{importantbox}

\subsection{本章小结}

最优推理策略依赖题目难度、模型能力和 verifier 质量。真正的 compute-optimal 系统需要在线估计边际收益，动态决定继续采样、修订、验证还是停止。

